\documentclass{situationist-preprint}

\title{The Thinking Animal Problem Solved}
\msid{PID-05-b}
\msfamily{PID-05}
\msdiscipline{Personal Identity and the Self}
\mstype{Companion manuscript}
\mskeywords{animalism; thinking animal; too many thinkers; constitution view; thinking parts; derivative thinking}

\begin{document}
\maketitle

\begin{abstract}
The thinking animal argument of PID-05-a relies on the premise that your present thoughts have at most one non-derivative thinker. Opponents of animalism reject it in two ways: by admitting two thinkers, a person and the animal that constitutes it, one of which thinks derivatively; or by locating the thinker in a part of the animal, the brain or cerebrum. We argue that both alternatives fail, and that the problem they are meant to solve, how a person and an animal can share thoughts without being two thinkers, has a simple solution once \enquote{person} is taken as a phase sortal of the organism: there are not two candidates. We also answer the worry that the organism's many large parts would be thinkers too.
\end{abstract}

\section{The problem}

There is a human organism where you are, and it thinks your thoughts (PID-05-a, Assumption~2.4). If you are not that organism, as psychological and constitution views hold, then two things think your thoughts. The person cannot know which of the two it is, and the animal, which is not a person on such views though it has the capacities of one, is mistaken whenever it thinks \enquote{I am a person} \citep{Olson1997,Snowdon1990}. This is the thinking animal problem. PID-05-a resolves it by Assumption~2.5: your thoughts have at most one non-derivative thinker; together with Assumptions~2.4 and~2.6 this yields Theorem~3.1, that you are an organism. The opponents of animalism deny Assumption~2.5 or deny that the thinker is the organism.

\section{Two thinkers, one derivative}

On the constitution view \citep{Baker2000}, the person is constituted by the animal without being identical with it, and the two share their thoughts; but the person has its first-person perspective non-derivatively, the animal only derivatively, by constituting something that has it. There are two thinkers, but only one in the primary sense.

\begin{proposition}\label{p:const}
The constitution view inverts the order of derivation.
\end{proposition}

\begin{proof}
A thing has a property derivatively if it has it in virtue of standing in a relation to another thing that has it non-derivatively. The animal's perceptual states, emotions and bodily feelings, which thoughts integrate, belong to it non-derivatively: it sees with its eyes and is afraid with its whole physiology, and the constitution theorist does not say that the animal sees only in virtue of constituting a seer. The first-person perspective is realized by the same brain processes, in the same life, as these states. Nothing in the animal's relation to its thoughts differs from its relation to its perceptions, and nothing in the person's relation to them is available to it except through the animal. If either has the thoughts derivatively, it is the person.
\end{proof}

A constituted person that thinks only derivatively is a thinker only in the sense in which a statue weighs something: in virtue of the stuff that does. Such a thing adds no second thinker to the count, but it adds nothing to the case for denying that you are the animal either. The thinker of your thoughts, in the only sense in which you think them, is the organism.

\section{Thinking parts}

A second alternative holds that the thinker is a part of the animal: the brain, or the cerebrum \citep{McMahan2002}. You would then be that part, would survive as a brain in a vat, and the animal would think only derivatively, as a man is said to be clever because his brain is.

PID-05-a replies that the brain thinks, if at all, as the hand grasps: it is the organ by which the organism thinks. We add an argument.

\begin{proposition}\label{p:parts}
A proper part of an organism that is not itself an organism is not a non-derivative thinker.
\end{proposition}

\begin{proof}
Thoughts integrate perceptual, affective and bodily states, and are states of a life: they draw on the organism's regulation of its own conditions, and they guide the organism's action as a whole (PID-05-a, Definition~2.1). A state of a life is non-derivatively a state of the thing whose life it is. A proper part that is not an organism has no life of its own; its activity is sustained by the organism's metabolism and has its function within it. So the thoughts it realizes are states of a life that is not its own, and it has them, if at all, derivatively.
\end{proof}

The same argument answers the worry that the animal's large proper parts, the animal minus a hand or minus a toe, would be thinkers too, on pain of indeterminacy as to which of very many overlapping things is you \citep{Unger1980}. Such parts are not organisms: they are not organized by a self-maintaining process of their own, and nothing is their life. By Proposition~\ref{p:parts}, none thinks non-derivatively. Only the organism does, and there is one organism where you are, since, by the composition principle of MET-06, parts compose one thing exactly when they share a common fate under the relevant dynamics, and only the whole organism has a fate of its own under the dynamics of life.

\section{The solution}

The thinking animal problem arises only for views on which you and the animal are two things that both satisfy the conditions for thinking your thoughts. On the account of PID-05-a, \enquote{person} is a phase sortal, true of the organism while it has the capacities for rational and self-conscious thought (Definition~2.2). The organism and the person are not two candidates for being you: the person is the organism in its personal phase, as a child is the human being in a phase of its life. There is no second thinker whose existence raises the question which of the two you are, and the animal's thought \enquote{I am a person} is true.

\begin{objection}
On the phase sortal view, the person of PID-01, whose persistence is given by psychological continuity, and the animal can come apart, as in the vat case (PID-05-a, Corollary~3.3). Then two things share your thoughts after all.
\end{objection}

\begin{reply}
They come apart only after the organism ceases to exist or ceases to think, so that at no time do two things think the same thoughts. Before the vat operation there is one thinker, the organism; afterwards there is a person in the vat, psychologically continuous with you, and no thinking organism. The person who continues in the vat is not you, since you are the organism (Theorem~3.2 of PID-05-a), though she has much of what matters (PID-06). Proposition~4.1 of PID-05-a states the reconciliation.
\end{reply}

\section{Conclusion}

The alternatives to the single-thinker premise either make the animal a derivative thinker, which inverts the order of derivation, or make a part of the animal the thinker, though a part without a life has no thoughts of its own. Once \enquote{person} is recognized as a phase sortal of the organism, there is only one thinker of your thoughts, and it is the animal you are.

\begin{thebibliography}{99}\small
\bibitem[Baker(2000)]{Baker2000} Baker, L.~R. (2000). \emph{Persons and Bodies: A Constitution View}. Cambridge University Press.
\bibitem[McMahan(2002)]{McMahan2002} McMahan, J. (2002). \emph{The Ethics of Killing: Problems at the Margins of Life}. Oxford University Press.
\bibitem[Olson(1997)]{Olson1997} Olson, E.~T. (1997). \emph{The Human Animal: Personal Identity without Psychology}. Oxford University Press.
\bibitem[Snowdon(1990)]{Snowdon1990} Snowdon, P.~F. (1990). Persons, animals, and ourselves. In C. Gill (ed.), \emph{The Person and the Human Mind}. Clarendon Press.
\bibitem[Unger(1980)]{Unger1980} Unger, P. (1980). The problem of the many. \emph{Midwest Studies in Philosophy}, 5, 411--467.
\end{thebibliography}

\end{document}
